i.e. makes @@M0@@ zero. Since the formation of @@M1@@ commutes with base change, we see that it is represented by the open subset @@M2@@ of @@M3@@. We may therefore suppose that @@M4@@ is already an epimorphism. But, now using the hypothesis that @@M5@@ is @@M6@@-flat, we find that the formation of @@M7@@ commutes with base change, and consequently the functor under consideration is represented by the open subset @@M8@@ of @@M9@@.

% Expose XIII local macros
\providecommand{\Ass}{\operatorname{Ass}}
\providecommand{\Hom}{\operatorname{Hom}}
\providecommand{\Ext}{\operatorname{Ext}}
\providecommand{\Proj}{\operatorname{Proj}}
\providecommand{\NS}{\operatorname{NS}}
\providecommand{\Tors}{\operatorname{Tors}}
\providecommand{\Picnum}{\operatorname{Pic}^{\mathrm{num}}}
\providecommand{\Piczero}{\operatorname{Pic}^{0}}
\providecommand{\Pictau}{\operatorname{Pic}^{\tau}}
\providecommand{\length}{\operatorname{length}}
\providecommand{\etale}{\text{\'etale}}

\newpage

